\begin{abstract}
% [Polish] Original: Electron density (ED) encodes the electronic-structure information that drives quantum molecular properties, yet computing it for every prediction is prohibitively expensive. We propose EDG, a cross-modal distillation framework that transfers ED knowledge from a pre-trained ED-aware teacher to geometry students as privileged information, requiring no ED computation at inference. Naive distillation, however, treats every teacher prediction as equally trustworthy and risks negative transfer; we therefore extend EDG to EDG++, which gates supervision through a pre-trained reliability estimator with adaptive thresholding, and prove (Proposition~1) a sufficient condition under which gating reduces an upper bound on gradient noise---yielding gains on 34/36 QM9 tasks, 55/60 rMD17 energy/force tasks, and 27--38\% tail-error reduction on the hardest 1\% of samples. As a stress test on the 166-atom mini-protein Chignolin---substantially out of the teacher's training distribution---naive EDG measurably degrades the baseline ($+\!4.0\%$ energy, $+\!2.8\%$ force MAE) while EDG++'s gating recovers a net improvement ($-\!6.8\%$ energy, $-\!2.9\%$ force), and the same frozen estimator additionally serves as an inference-time OOD detector during MD rollouts.
Electron density (ED), which describes the spatial distribution of electrons in a molecule, is crucial for accurately predicting the energies and forces in molecular force fields (MFF). However, existing MLFF methods exploit only atom-level geometric quantities and ignore the electronic-level information that governs these properties. To bridge this gap, we propose EDG, a cross-modal knowledge distillation framework that transfers ED knowledge from a pre-trained ED-aware teacher to geometry-based students at training time; at inference, only the geometry student is used, avoiding the prohibitive cost of quantum chemical calculations. Since naive distillation treats every teacher prediction as equally trustworthy and can lead to negative transfer, we further propose EDG++, which gates the distillation signal through a pre-trained reliability estimator with adaptive thresholding; a theoretical analysis shows that the gating reduces an upper bound on the gradient noise. Extensive experiments on QM9 and rMD17 show that the framework can be directly integrated into a wide range of geometry-based models (SchNet, Equiformer, SphereNet, ViSNet), significantly improving their geometric representation learning capability with a maximum average improvement of 33.7\%, with consistent gains on 34/36 QM9 task-model combinations and 55/60 rMD17 energy/force tasks. On the 166-atom mini-protein Chignolin, far outside the teacher's training distribution, EDG++ further demonstrates the rescue effect of reliability gating against the negative transfer induced by naive distillation. EDG/EDG++ offers a general framework for injecting electronic-level inductive bias into geometric molecular models, with applications to molecular property prediction, ML force field learning, and biomolecular dynamics simulation.
\end{abstract}

\begin{IEEEkeywords}
Electron density, molecular geometry learning, knowledge distillation, reliability-aware learning, machine learning force fields.
\end{IEEEkeywords}
